\subsection{Qué significa la brecha de 0.35}

El resultado central no es que el modelo generalice mal, que ya lo documentaron
Beery et al.~\cite{beery2018recognition}, sino que ahora se puede decir cuánto
de esa pérdida es realmente por el cambio de ubicación: el control de
distribución descarta la explicación alternativa, porque solo $0.033$ de los
$0.3485$ se explican por el cambio de prevalencia.

El contraste con \texttt{car} es el argumento más fuerte. Si la caída viniera de
un cambio general de dominio (otra cámara, otra iluminación, otro encuadre), el
auto debería caer también. No cae. La diferencia no es el dominio: el auto es
rígido, grande y de alto contraste, mientras que un \texttt{squirrel} a veinte
metros es una mancha que se distingue porque \emph{esa} mancha, en \emph{ese}
tronco, históricamente fue una ardilla.

Grad-CAM lo confirma de forma independiente, y conviene decir hasta dónde. Una
razón de masa mediana de $1.00$ en ubicaciones nuevas no prueba que el modelo
mire el fondo: prueba que su evidencia no está más concentrada sobre el animal
que si estuviera repartida al azar. Es más débil que ``mira el fondo'', pero es
lo que el dato sostiene; para animales chicos ($0.89$) sí cruza el umbral de la
anti-localización.

\subsection{Lo que la métrica agregada no puede expresar}

El caso de \texttt{deer} en \texttt{trans\_test} es instructivo más allá de este
trabajo: cero cajas anotadas y 646 detecciones emitidas. La AP no está definida y
se reporta como tal, pero ningún valor de AP (ni cero, ni indefinido) comunica
que el modelo inventa ciervos. Con \texttt{bird} pasa el caso inverso: 17\,086
detecciones para 180 cajas y AP $= 0.0132$, que un lector entiende como ``no
detecta pájaros'' cuando el problema es el exceso de detecciones. De ahí que la
tabla por especie incluya las detecciones junto al soporte: convierte un AP bajo
en un diagnóstico, porque con pocas detecciones es falta de recall y con muchas,
de precisión, y las dos se arreglan distinto.

\subsection{Decisiones que podrían haber salido de otra forma}

\emph{La resolución.} 640 px se tomó porque es la de los pesos preentrenados y
permite un lote de 16. El barrido por tamaño muestra su costo; el experimento que
lo cerraría, reentrenar a 960 px, no se hizo, así que se afirma el diagnóstico y
no la mitigación.

\emph{El tramo grande rinde peor que el mediano}, en las dos particiones
($0.7341$ contra $0.8033$, y $0.2852$ contra $0.4555$), y no era lo esperado. La
hipótesis plausible es que un animal que ocupa más del 5\% del cuadro está muy
cerca del lente: entra parcialmente fuera de cuadro, con desenfoque de movimiento
y, de noche, sobreexpuesto por el infrarrojo. Cuál domina no se midió, así que
queda como observación.

\emph{Una predicción equivocada sobre la oclusión.} Se diseñó el modo
\texttt{inferior} esperando que fuera el caso más duro, porque rompe la silueta
además de tapar la textura. Medido, al 25\% es claramente más benigno que tapar
el centro ($0.3699$ contra $0.2375$): el modelo se apoya más en la textura
interior que en el contorno, y al 50\% la diferencia se desvanece. Se registra
porque la hipótesis estaba escrita antes de medir.

\emph{Los negativos explícitos.} Las imágenes vacías podrían sumarse como
fondo; no se hizo porque habría cambiado a la vez el tamaño del conjunto y la
proporción fondo/objeto. Dado el volumen de falsos positivos en terreno nuevo,
es la ablación pendiente más prometedora.

\subsection{Límites de lo que se midió}

\emph{Una sola corrida de entrenamiento.} Los intervalos de confianza
capturan la variabilidad del \emph{muestreo de imágenes}, no la de la
\emph{inicialización}: no se entrenaron varias semillas. La brecha de $0.35$ es
mucho mayor que cualquier variación por semilla plausible, pero las diferencias
chicas de este informe, como el tramo grande contra el mediano, no se pueden
atribuir con esa evidencia.

\emph{La oclusión es sintética.} Un rectángulo gris uniforme es más fácil de
ignorar, o de confundir, que una rama real con textura. La dirección del
resultado es contundente; su magnitud exacta no se traslada a oclusión natural.

\emph{La ocupación del módulo ecológico es naive} y debe leerse como cota
inferior. El error del reporte se midió sobre las dos particiones: en
\texttt{cis\_test}, el caso favorable, va de $+2\%$ a $+100\%$ y siempre
sobreestima; en \texttt{trans\_test} es mayor, hasta $+281\%$ en \texttt{bird}, y
aparece en las dos direcciones, porque el modelo inventa eventos de las especies
comunes y a la vez pierde los de \texttt{badger} ($-75\%$), la rara.
